\chapter{Conclusão}

Este trabalho apresenta uma proposta de sistema de tradução automática voltado para documentos jurídicos, utilizando técnicas de Processamento de Linguagem Natural (NLP) para capturar nuances culturais e legais entre o português e o inglês. A partir da análise do contexto jurídico e dos desafios da tradução, identifica-se que muitos termos e expressões legais requerem uma interpretação cuidadosa que vai além de uma tradução literal, reforçando a relevância de métodos avançados como modelos de aprendizado profundo.

O desenvolvimento do sistema inclui etapas de coleta e pré-processamento de dados, modelagem do modelo de tradução e avaliação da eficácia do sistema. Com isso, busca-se criar uma solução que auxilie na interpretação e disseminação de documentos jurídicos brasileiros em um contexto global, tornando essas informações mais acessíveis e compreensíveis para um público diversificado.

Os resultados, tanto quantitativos (como a métrica BLEU) quanto qualitativos (a partir da avaliação de profissionais da área jurídica), indicam que o sistema realiza traduções mais adequadas e contextualizadas em comparação com ferramentas de tradução convencionais. A adesão a um banco de dados ontologicamente estruturado mostra-se essencial para a compreensão semântica, proporcionando uma base sólida para as traduções jurídicas automatizadas.

Por fim, este trabalho contribui para a área de tradução automática ao explorar uma aplicação específica e complexa, como o contexto jurídico, sugerindo que o uso de tecnologias baseadas em NLP aprimora a precisão e a relevância das traduções em documentos especializados. Futuras pesquisas podem explorar a inclusão de outros idiomas e o refinamento dos modelos, além de investigar novas métricas de avaliação para capturar aspectos mais subjetivos presentes nas traduções legais.

